\documentclass[11pt]{scrartcl}
\usepackage[spanish]{babel}
\usepackage[sexy, fancy]{evan}
\usepackage{pgfplots}
\usepackage{float}
\pgfplotsset{compat=1.15}
\usepackage{mathrsfs}
\usetikzlibrary{arrows}
\usepackage{graphics}
\usepackage{csquotes} %Edición de textos
\usepackage{multicol} %multicolumnas
\usepackage{tikz}
\usepackage[dvipsnames]{xcolor}
\definecolor{red1}{RGB}{255, 153, 153}
\definecolor{green1}{RGB}{204, 255, 204}
\definecolor{blue1}{RGB}{204, 255, 255}
\definecolor{yellow1}{RGB}{255, 247, 160}

\definecolor{red2}{RGB}{255, 102, 102}
\definecolor{green2}{RGB}{108, 255, 108}
\definecolor{blue2}{RGB}{94, 204, 255}
\definecolor{yellow2}{RGB}{255, 250, 104}

\definecolor{red2.5}{RGB}{255,76,76}
\definecolor{green2.5}{RGB}{54, 247, 54}
\definecolor{blue2.5}{RGB}{51, 189, 255}
\definecolor{yellow2.5}{RGB}{255, 242, 52}

\definecolor{red3}{RGB}{255, 51, 51}
\definecolor{green3}{RGB}{0, 240, 0}
\definecolor{blue3}{RGB}{9, 175, 255}
\definecolor{yellow3}{RGB}{255, 234, 0}

\definecolor{red3.5}{RGB}{229, 25, 25}
\definecolor{green3.5}{RGB}{0, 194, 0}
\definecolor{blue3.5}{RGB}{4, 143, 209}
\definecolor{yellow3.5}{RGB}{255,220,0}

\definecolor{red4}{RGB}{204, 0, 0}
\definecolor{green4}{RGB}{0, 149, 0}
\definecolor{blue4}{RGB}{0, 111, 164}
\definecolor{yellow4}{RGB}{255, 206, 0}

\usepackage{graphicx} % Required for inserting images

\usepackage{amsmath} %Para poner acentos

\usepackage{amssymb} %Geogebra
\usepackage{asymptote}
\usepackage{amsthm}
\usepackage{chngcntr} %Descripciones de figuras
\title{\texttt{Tarea 10
\\ Cálculo Diferencial e Integral II}}
\author{Luis David Uicab Martínez. 
\\Lic. en Matemáticas.}
\date{02 de mayo del 2026}

\begin{document}
\maketitle
\textbf{Lema I}. Dado $N\in \NN$ arbitrario, y $(n_k)$ una sucesión de números naturales
estrictamente creciente, existe $k' \in \NN$ tal que para todo $k > k'$, $n_k > N$. \\ \\
\textbf{Demostración.} Procedamos por contradicción
y supongamos que $n_{k}\leq N$ para todo $k\in \NN$. Así, $n_k$ puede ser uno de $n$ valores distintos. Luego, tomando los primeros $N+1$ términos de $(n_k)$,
podemos afirmar que al menos dos tienen el mismo valor. Pero esto es una contradicción, puesto que los $n_k$ son por definición estrictamente
crecientes, y en particular, distintos entre sí. En consecuencia, existe $k'$ tal que $n_{k'}>N$. Puesto que como ya mencionamos, los $n_k$'s son crecientes,
decimos que para toda $k > k'$, $n_k > N$. $\blacksquare$
\begin{mdframed}[style=mdpurplebox, frametitle={Ejercicio 01, Tarea 10, Cálculo II}]
    (\textbf{Identidades del ángulo medio}) Prueba lo siguiente:
    \begin{enumerate}[(a)]
        \item  $\sin\left(\frac{x}{2}\right)=\pm\sqrt{\frac{1-\cos\left(x\right)}{2}}.$
        \item  $\cos\left(\frac{x}{2}\right)=\pm\sqrt{\frac{1+\cos\left(x\right)}{2}}.$
        \item  $\tan\left(\frac{x}{2}\right)=\pm\sqrt{\frac{1-\cos\left(x\right)}{1+\cos\left(x\right)}}.$
    \end{enumerate}
\end{mdframed}
\textbf{Demostración.}
\begin{enumerate}[(a)]
    \item Del inciso (a) del corolario (3) sabemos que $\sin^2(\alpha)=\left(1-\cos(2\alpha)\right)/2$. Así,
    tomando $\alpha=x/2$ obtenemos que
    \[
    \sin^2\left(\frac{x}{2}\right)=\frac{1-\cos(x)}{2} \iff \sin\left(\frac{x}{2}\right)=\pm \sqrt{\frac{1-\cos(x)}{2}}.
    \]
    \item Del inciso (b) del corolario (3) decimos que $\cos^2(\alpha)=\left(1+\cos(2\alpha)\right)/2$. Luego,
    basta tomar $\alpha=x/2$ para afirmar que
    \[
    \cos^2\left(\frac{x}{2}\right)=\frac{1+\cos(x)}{2} \iff \cos\left(\frac{x}{2}\right)=\pm \sqrt{\frac{1+\cos(x)}{2}}.
    \]
    \item Del inciso (c) del corolario (3), tenemos la igualdad $\tan^2(\alpha)=\left(1-\cos(2\alpha)\right)/\left(1+\cos(2\alpha)\right)$.
    Nuevamente, tomando $\alpha=x/2$ se obtiene que
    \[
    \tan^2\left(\frac{x}{2}\right)=\frac{1-\cos(x)}{1+\cos(x)} \iff \tan\left(\frac{x}{2}\right)=\sqrt{\frac{1-\cos(x)}{1+\cos(x)}},
    \]
    con lo que damos por concluida la demostración. $\blacksquare$    
\end{enumerate}

\begin{mdframed}[style=mdpurplebox, frametitle={Ejercicio 02, Tarea 10, Cálculo II}]
    Considera una sucesión $\left(a_{n}\right).$ Prueba que si $a_{n}\rightarrow\pm\infty$
    y $\left(a_{n_{k}}\right)$ es una subsucesión de $\left(a_{n}\right),$
    entonces $a_{n_{k}}\rightarrow\text{\ensuremath{\pm}}\infty.$
\end{mdframed}
\textbf{Demostración.} Supongamos que $a_n \rightarrow \infty$. Así, para $M>0$
arbitrario, existe $N\in \NN$ tal que para todo $n > N$
\[
a_n > M. \tag{1}
\]
Por otra parte, sea $(a_{n_k})$ subsucesión de $(a_n)$. A proiri, por el lema I existe $k' \in \NN$ tal que para toda $k > k'$, $n_{k'}>N$.  Luego, como (1)
se cumple para todo $n > N$, de las desigualdades mencionadas con anterioridad, afirmamos que para todo $k > k'$
\[
a_{n_k} > M,
\]
con lo que concluimos, $(a_{n_k}) \rightarrow \infty$. La prueba de que si $(a_n)\rightarrow -\infty$, $(a_{n_k})\rightarrow -\infty$, es similar
(en tal caso se toma $M < 0$ arbitrario, y afirmamos que para todo $n > N$, $a_n < M$; el resto es análogo). $\blacksquare$
\newpage
\begin{mdframed}[style=mdpurplebox, frametitle={Ejercicio 03, Tarea 10, Cálculo II}]
    \textbf{(Agrupamiento de una serie)} Dada una serie $\sum_{n=1}^{\infty}a_{n}$
    y una sucesión creciente de números naturales $n_{1}<n_{2}<n_{3}<...,$
    definamos 
    \begin{align*}
        A_{1} & =a_{1}+\cdots a_{n_{1}},\\
        A_{2} & =a_{n_{1}+1}+\cdots+a_{n_{2}},\\
        A_{3} & =a_{n_{2}+1}+\cdots+a_{n_{3}},\\
        & \vdots\\
        A_{k} & =a_{n_{k-1}+1}+\cdots+a_{n_{k}}.
    \end{align*}
    A la serie $\sum_{k=1}^{\infty}A_{k}$ le llamaremos un \textit{agrupamiento}\index{agrupamiento de una serie}
    de la serie $\sum_{n=1}^{\infty}a_{n}.$ El porqué de la terminología
    \enquote{agrupamiento de la serie$\sum_{n=1}^{\infty}a_{n}$} queda claro
    si denotamos a la serie $\sum_{k=1}^{\infty}A_{k}$ como 
    \[
    \left(a_{1}+\cdots+a_{n_{1}}\right)+\left(a_{n_{1}+1}+\cdots+a_{n_{2}}\right)+\left(a_{n_{2}+1}+\cdots+a_{n_{3}}\right)+\cdots
    \]
    Prueba que si la serie $\sum_{n=1}^{\infty}a_{n}$ converge a $L,$
    entonces todos sus agrupamientos también convergen a $L.$
\end{mdframed}
\textbf{Demostración.} Supongamos que $\sum_{n=1}^{\infty}a_n = L$. Así, dado un $\varepsilon > 0$ arbitrario, existe $N\in \NN$, tal que
para todo $s > N$,
\[
\left|\sum_{n=1}^{s}a_n - L\right| < \varepsilon. \tag{1}
\]
Luego, por el lema I, existe $M \in \NN$ tal que para todo $m > M$, $n_m > N$. En consecuencia, añadiendo el hecho de que
(1) se cumple para todo $s > N$ y de la definición de los $A_k$, afirmamos que, para todo $m > M$,
\begin{align*}
    \left|\sum_{n=1}^{n_m}a_n - L\right| < \varepsilon &\implies \left|\left(\sum_{n=1}^{n_1}a_n + \sum_{n=n_1+1}^{n_2}a_n + ... + \sum_{n=n_{m-1}+1}^{n_m}a_n\right) - L\right| < \varepsilon \\
    &\implies \left|\left(\sum_{k=1}^{m}A_k\right) - L\right| < \varepsilon,
\end{align*}
con lo que concluimos,
\[
\sum_{k=1}^{\infty}A_k = L. \quad \blacksquare
\]
\newpage
\begin{mdframed}[style=mdpurplebox, frametitle={Ejercicio 04, Tarea 10, Cálculo II}]
    Determina la convergencia de las series siguientes: 
    \begin{multicols}{2}
        \begin{enumerate}[(a)]
            \item  $\sum_{n=1}^{\infty}$$\frac{n}{n+1}$.
            \item  $\sum_{n=1}^{\infty}\frac{2^{n}+3^{n}}{4^{n}}.$
            \item  $\sum_{n=1}^{\infty}\frac{5n}{n^{2}+1}.$
            \item  $\sum_{n=2}^{\infty}\frac{\text{cos}\left(n^{3}+9\right)}{n^{3}+9}.$
        \end{enumerate}
    \end{multicols}
\end{mdframed}
\textbf{Demostración.}
\begin{enumerate}[(a)]
    \item Definamos las sucesiones $(a_n)=1$ y $(b_n)=1/(n+1)$. Puesto que $(a_n)$ es la constante 1,
    se tiene que $(a_n)\rightarrow 1$. Por otra parte puesto que $(b_n)$ es una sucesión geométrica, se sigue que converge
    a 0. De lo anterior, por el álgebra de límites para sucesiones, decimos que la sucesión $(a_n-b_n)=(1-1/(n+1))=(n/(n+1))$
    converge y además
    \[
    (a_n-b_n) \rightarrow 1.
    \]
    Luego, de la afirmación anterior, hacemos notar que el $n$-ésimo término de la serie $\sum_{n=1}^{\infty}(n/(n+1))$ es distinto de 0,
    con lo que, por la contrapositiva de la proposición 73, dicha serie no converge.
    \item En principio, notemos que
    \[
    \sum_{n=0}^{\infty}\left(\frac{1}{2}\right)^n, \quad \sum_{n=0}^{\infty}\left(\frac{3}{4}\right)^n,
    \]
    son series de la forma $\sum_{n=1}^{\infty}x^n$ tales que $|x|<1$ (es decir, $|1/2|,|3/4|<1$), y en consecuencia, por la proposición 74,
    convergen. Luego, por el ejercicio 7 de la tarea 9, esto implica que
    \[
    \sum_{n=1}^{\infty}\left(\frac{1}{2}\right)^n \quad \text{y} \quad \sum_{n=1}^{\infty}\left(\frac{3}{4}\right)^n
    \]
    convergen. Así, puesto que
    \begin{align*}
        \left(\frac{1}{2}\right)^n + \left(\frac{3}{4}\right)^n &= \left(\frac{2}{4}\right)^n+\left(\frac{3}{4}\right)^n \\
        &= \frac{2^n}{4^n}+\frac{3^n}{4^n} \\
        &= \frac{2^n+3^n}{4^n},
    \end{align*}
    concluimos que
    \[
    \sum_{n=1}^{\infty}\left(\frac{2^n+3^n}{4^n}\right)=\sum_{n=1}^{\infty}\left(\frac{1}{2}\right)^n + \sum_{n=1}^{\infty}\left(\frac{3}{4}\right)^n,
    \]
    y en adición, al ser la suma de dos series convergentes, ella converge.
    \item Puesto que $0 \leq n^2+1 \leq n^2 +n$ para todo $n \in \NN$, se tiene que
    \[
    0 \leq \frac{5n}{n^2+n} \leq \frac{5n}{n^2+1}, \quad \forall n \in \NN. \tag{1}
    \]
    Por otra parte, notemos que
    \[
    \sum_{n=1}^{\infty}\frac{5n}{n^2+n}=5\sum_{n=1}^{\infty}\frac{n}{n(n+1)}=5\sum_{n=1}^{\infty}\frac{1}{n+1}
    \]
    es una serie armónica, y en consecuencia, diverge, con lo que, dado el criterio de comparación directo y la igualdad en (1) se sigue que
    \[
    \sum_{n=1}^{\infty}\frac{5n}{n^2+1}
    \]
    diverge.
    \item A priori, hacemos ver que $|\cos(z)|\leq 1$. Por otra parte, notemos que $0 \leq 1/(n^3+9) \leq 1/n^3, n\in \NN$.
    Luego, de ambas desigualdades, se tiene que
    \[
    0 \leq \frac{\left|\cos(n^3 + 9)\right|}{n^3+9} \leq \frac{1}{n^3}, \forall n \in \NN. \tag{2}
    \]
    Luego, puesto que por la proposición 77
    \[
    \sum_{n=1}^{\infty}\frac{1}{n^3}
    \]
    converge y dado que por el ejercicio 7 de la tarea 9 esto implica que dicha serie converge desde $n=2$, del criterio de comparación directo y de la
    desigualdad en (2) concluimos que
    \[
    \sum_{n=2}^{\infty}\frac{\left|\cos(n^3 + 9)\right|}{n^3+9} = \sum_{n=2}^{\infty}\left|\frac{\cos(n^3 + 9)}{n^3+9}\right|
    \]
    también converge. Luego, como de lo anterior, $\sum_{n=2}^{\infty}\frac{\cos(n^3 + 9)}{n^3+9}$ converge absolutamente, por la proposición 84 converge.
    $\blacksquare$
\end{enumerate}
\begin{mdframed}[style=mdpurplebox, frametitle={Ejercicio 05, Tarea 10, Cálculo II}]
    \textbf{(Expansión decimal)} 
    \begin{enumerate}[(a)]
        \item Prueba que si $\left(a_{n}\right)$ es una sucesión de números enteros
        con $0\leq a_{n}\leq 9,$ entonces la serie 
        \[ 
        \sum_{n=1}^{\infty}\frac{a_{n}}{10^{n}}
        \]
        converge a un número del intervalo $\left[0,1\right].$ 
        \item Prueba que si $x\in\left[0,1\right],$ entonces existe una sucesión
        de números enteros $\left(a_{n}\right)$ con $0\leq a_{n}\leq9$ tal que 
        \[
        x=\sum_{n=1}^{\infty}\frac{a_{n}}{10^{n}}.
        \]
        \item Prueba que si $x\geq0,$ entonces existen un entero no positivo
        $N$ y una sucesión de números enteros $\left(a_{n}\right)_{n=N}^{\infty}$
        con $0\leq a_{n}\leq9$ tales que 
        \begin{equation}
            x=\sum_{n=N}^{\infty}\frac{a_{n}}{10^{n}}.\label{Ec: Expansi=0000F3n decimal}
        \end{equation}
        La serie de (\ref{Ec: Expansi=0000F3n decimal}) es la llamada \textit{expansión
        decimal}\index{expansión decimal de un número} del número $x$ y
        se denota por $a_{N}a_{N+1}...a_{0}.a_{1}a_{2}a_{3}...$ La expansión
        decimal de un número negativo $x$ se define como el negativo de la
        expansión decimal del número $-x.$
    \end{enumerate}
    \textbf{Sugerencia:} Para probar el inciso (b), toma $a_{1}$ como
    el mayor entero tal que $0\leq a_{1}\leq 9$ y 
    \[
    \frac{a_{1}}{10}\leq x.
    \]
    Luego toma $a_{2}$ como el mayor entero tal que $0\leq a_{2}\leq9$
    y 
    \[
    \frac{a_{1}}{10}+\frac{a_{2}}{10^{2}}\leq x,
    \]
    etc. Para probar el inciso (c) usa el inciso (b). Recuerda el ejercicio
    71 de las notas de Cálculo I.    
\end{mdframed}
\textbf{Demostración.}
\begin{enumerate}[(a)]
    \item A priori, puesto que $0\leq a_n \leq 9$, notemos que
    \begin{equation*}
        0 \leq \frac{a_n}{10^n} \leq \frac{9}{10^n}, \forall n \in \NN. \tag{1}
    \end{equation*}
    Por otra parte, puesto que $(1/10) < 1$, se sigue que la serie geométrica $\sum_{n=0}^{\infty}\frac{1}{10^n}$ converge
    y además, del ejercicio 7 de la tarea 9 y de la proposición 74
    \begin{align*}
        \frac{10}{9} &= \frac{1}{1-\frac{1}{10}} \\
        &= \sum_{n=0}^{\infty}\frac{1}{10^n} \\
        &= 1 + \sum_{n=1}^{\infty}\frac{1}{10^n}.
    \end{align*}
    De dicha igualdad, afirmamos que
    \begin{align*}
        \sum_{n=1}^{\infty}\frac{9}{10^n} &= 9\cdot \sum_{n=1}^{\infty}\frac{1}{10^n} \\
        &= 9\left(\frac{10}{9}-1\right) \\
        &= 1. \tag{2}
    \end{align*}
    Puesto que de lo anterior, sabemos que $\sum_{n=1}^{\infty}\frac{9}{10^n}$ converge, por el criterio de comparación directa
    y la desigualdad en (1), se sigue que $\sum_{n=1}^{\infty}\frac{a_n}{10^n}$ también converge. Por otra parte,
    notemos que
    \begin{equation*}
        0 \leq \sum_{n=1}^{m}\frac{a_n}{10^n} \leq \sum_{n=1}^{m}\frac{9}{10^n}, \forall n \in \NN.
    \end{equation*}
    Así, puesto que ambas sucesiones convergen (puesto que son las sucesiones asociadas a las sumas parciales de las series antes
    mencionadas), de (2) y del ejercicio 2 de la tarea 9, decimos que
    \begin{equation*}
        0 \leq \sum_{n=1}^{\infty}\frac{a_n}{10^n} \leq \sum_{n=1}^{\infty}\frac{9}{10^n} = 1,
    \end{equation*}
    con lo que concluimos, $\sum_{n=1}^{\infty}\frac{a_n}{10^n} \in [0,1]$.
    \item Sea $\varepsilon > 0$ arbitrario. Existe $N \in \NN$ tal que para todo $m > N$, se cumple que
    \begin{equation*}
        \left|\sum_{n=1}^{m}\frac{a_n}{10^n} - x\right| < \varepsilon
    \end{equation*}
\end{enumerate}
\end{document}